\section*{第 \ref{ch:g7:fractions} 章——分数：比较与加法}

\begin{solution}{exo:g7:fractions:1}
$\dfrac{2}{5} = \dfrac{8}{20}$; \quad
$\dfrac{18}{24} = \dfrac{3}{4}$; \quad
$\dfrac{7}{3} = \dfrac{28}{12}$; \quad
$\dfrac{5}{9} = \dfrac{20}{36}$.
\end{solution}

\begin{solution}{exo:g7:fractions:2}
$\dfrac{12}{16} = \dfrac34$; \quad
$\dfrac{30}{45} = \dfrac23$; \quad
$\dfrac{27}{36} = \dfrac34$; \quad
$\dfrac{48}{12} = 4$.
\end{solution}

\begin{solution}{exo:g7:fractions:3}
$\dfrac{15}{35} = \dfrac37$ 和 $\dfrac{21}{49} = \dfrac37$：相等。

$\dfrac{16}{28} = \dfrac47$ 和 $\dfrac{20}{36} = \dfrac59$：比较时，
$\frac47 = \frac{36}{63}$ 与 $\frac59 = \frac{35}{63}$——
不相等。
\end{solution}

\begin{solution}{exo:g7:fractions:4}
$\dfrac{7}{11} < \dfrac{9}{11}$（\omterm{def:g4:fractions:def}{分母}相同）。

$\dfrac56 = \dfrac{15}{18} > \dfrac{11}{18}$.

$\dfrac{13}{12} > 1$，而 $\dfrac{19}{20} < 1$，所以
$\dfrac{13}{12} > \dfrac{19}{20}$.
\end{solution}

\begin{solution}{exo:g7:fractions:5}
$\dfrac59 + \dfrac29 = \dfrac79$; \quad
$\dfrac{11}{7} - \dfrac47 = \dfrac77 = 1$; \quad
$\dfrac58 + \dfrac78 = \dfrac{12}{8} = \dfrac32$.
\end{solution}

\begin{solution}{exo:g7:fractions:6}
$\dfrac23 + \dfrac56 = \dfrac46 + \dfrac56 = \dfrac96 = \dfrac32$.

$\dfrac{7}{10} - \dfrac25 = \dfrac{7}{10} - \dfrac{4}{10} =
\dfrac{3}{10}$.

$\dfrac34 + \dfrac{5}{12} = \dfrac{9}{12} + \dfrac{5}{12} =
\dfrac{14}{12} = \dfrac76$.
\end{solution}

\begin{solution}{exo:g7:fractions:7}
$3 - \dfrac54 = \dfrac{12}{4} - \dfrac54 = \dfrac74$; \quad
$1 + \dfrac38 = \dfrac{11}{8}$; \quad
$2 - \dfrac76 = \dfrac{12}{6} - \dfrac76 = \dfrac56$.
\end{solution}

\begin{solution}{exo:g7:fractions:8}
$5 \times \dfrac{3}{10} = \dfrac{15}{10} = \dfrac32$; \quad
$\dfrac78 \times 4 = \dfrac{28}{8} = \dfrac72$; \quad
$\dfrac23$ 的 $45$ 等于 $\dfrac{2 \times 45}{3} = \dfrac{90}{3} = 30$。
\end{solution}

\begin{solution}{exo:g7:fractions:9}
两者合计：$\dfrac14 + \dfrac38 = \dfrac28 + \dfrac38 = \dfrac58$ 的
披萨。剩余：$1 - \dfrac58 = \dfrac38$。
\end{solution}

\begin{solution}{exo:g7:fractions:10}
倒出：$2 \times \dfrac16 = \dfrac26 = \dfrac13$ L。剩余：
\[
\frac34 - \frac13 = \frac{9}{12} - \frac{4}{12} = \frac{5}{12}
\text{ L}.
\]
\end{solution}

\begin{solution}{exo:g7:fractions:11}
将正数 $\frac12$ 加到 $\frac35$ 上，结果必须比 $\frac35$
\emph{更大}。但 $\frac47 < \frac35$（事实上
$\frac47 = \frac{20}{35}$ 与 $\frac35 = \frac{21}{35}$）：所给的和
竟小于其中一个加数——不可能。
\end{solution}

\begin{solution}{exo:g7:fractions:12}
逐步计算：
\[
\frac12 + \frac14 = \frac34, \qquad
\frac34 + \frac18 = \frac68 + \frac18 = \frac78, \qquad
\frac78 + \frac1{16} = \frac{14}{16} + \frac1{16} = \frac{15}{16}.
\]
部分和 $\frac12, \frac34, \frac78, \frac{15}{16}, \dots$ 依次，
每次都向 $1$ 靠近剩余距离的\omterm{def:g3:division:half}{一半}：每一步中，缺少的部分
被截去\omterm{def:g3:division:half}{一半}。这些和无限接近 $1$，却永远达不到它。
\end{solution}

\begin{solution}{pb:g7:fractions:1}
\textbf{1.} (a) $\frac12 + \frac14 + \frac14 = \frac24 + \frac14
+ \frac14 = \frac44 = 1$：符合要求。(b) $\frac14 + \frac18 +
\frac18 + \frac12 = \frac28 + \frac18 + \frac18 + \frac48 =
\frac88 = 1$：符合要求。

\textbf{2.} $\frac12 + \frac14 + \frac18 = \frac48 + \frac28 +
\frac18 = \frac78$：缺少一个八分音符。

\textbf{3.} 一小节为 $1 = \frac{16}{16}$：十六个十六分
音符。八分音符为 $\frac18 = \frac{2}{16}$：两个十六分音符。

\textbf{4.} 附点二分音符：$\frac12 + \frac14 = \frac34$。附点
四分音符：$\frac14 + \frac18 = \frac28 + \frac18 = \frac38$。

\textbf{5.} 附点二分音符的时值为 $\frac34$（第 4 题）：
恰好是一个华尔兹小节。其他填法，例如：四分音符
$+$ 四分音符 $+$ 四分音符（$\frac14 + \frac14 + \frac14 =
\frac34$），或四分音符 $+$ 八分音符 $+$ 八分音符 $+$ 四分音符
($\frac14 + \frac18 + \frac18 + \frac14 = \frac{2 + 1 + 1 +
2}{8} = \frac68 = \frac34$).

\textbf{6.} $\frac38 + \frac18 + \frac14 = \frac38 + \frac18 +
\frac28 = \frac68 = \frac34$：缺少
一个四分音符（$\frac14$）。

\textbf{7.} 附点八分音符：$\frac18 + \frac{1}{16} =
\frac{2}{16} + \frac{1}{16} = \frac{3}{16}$。四分音符为
$\frac{4}{16}$，更长。

\textbf{8.} $\frac12 + \frac18 = \frac48 + \frac18 = \frac58$,
即 $\frac{10}{16}$：十个十六分音符。

\textbf{9.} $1 - \frac{11}{16} = \frac{16}{16} - \frac{11}{16}
= \frac{5}{16}$，即小节中的休止部分。

\textbf{10.} 两个例子合计 $\frac{2 + 1 + 1}{16} =
\frac{4}{16}$：合法。所有填满四个十六分之一的节奏，
均由 $2$ 与 $1$ 两种片段组成：
\[
2{+}2, \quad 2{+}1{+}1, \quad 1{+}2{+}1, \quad 1{+}1{+}2, \quad
1{+}1{+}1{+}1 :
\]
共五种节奏。

\textbf{11.} 填满 $n$ 个十六分之一的节奏，末尾为
一个十六分音符——其前部分填满 $n - 1$ 个——或为
一个八分音符——其前部分填满 $n - 2$ 个。每种节奏恰好被计数一次，
于是得
$\text{种数}(n) = \text{种数}(n-1) + \text{种数}(n-2)$。列表如下：
\begin{center}
\renewcommand{\arraystretch}{1.2}
\begin{tabular}{c|cccccccc}
$n$ & $1$ & $2$ & $3$ & $4$ & $5$ & $6$ & $7$ & $8$ \\
\hline
种数 & $1$ & $2$ & $3$ & $5$ & $8$ & $13$ & $21$ & $34$ \\
\end{tabular}
\end{center}
半音符（$8$ 个十六分之一）共有 $34$ 种敲法。
（$n = 4$ 时恰好得到第 10 题中的五种节奏。）

\textbf{12.} 每个数都是前两个数之和：
$3 + 5 = 8$, $5 + 8 = 13$, $8 + 13 = 21$, $13 + 21 = 34$——
第 11 题所证明的规律，已有千年之久，
如今被称为斐波那契规律。

\textbf{13.} $\frac16 + \frac{1}{12} = \frac{2}{12} +
\frac{1}{12} = \frac{3}{12} = \frac14$：与
两个平直八分音符（$\frac18 + \frac18 = \frac14$）时值相同。
摇摆式互换是合法的。

\textbf{14.} 每个三连音音符的时值为 $\frac{1}{12}$：验算如下，
$3 \times \frac{1}{12} = \frac{3}{12} = \frac14$
（\cref{prop:g7:fractions:mult}）。对二分音符而言，每个
音符时值为 $\frac16$，因为 $3 \times \frac16 = \frac36 = \frac12$。

\textbf{15.} 附点四分音符 $+$ 四分音符 $+$ 四分音符：以八分音符计，
$3 + 2 + 2 = 7$ 个八分音符 $= \frac78$：合法。二分音符与四分音符
分别相当于 $4$ 个与 $2$ 个八分音符——两者都是\emph{\omterm{def:g3:numbers:evenodd}{偶数}}。
任意\omterm{def:g3:numbers:evenodd}{偶数}之和仍为\omterm{def:g3:numbers:evenodd}{偶数}，但七八拍小节需要
$7$ 个八分音符，是\omterm{def:g3:numbers:evenodd}{奇数}，不可能。要填满\omterm{def:g3:numbers:evenodd}{奇数}个
八分音符的小节，必须出现\omterm{def:g3:numbers:evenodd}{奇数}时值的音符——八分音符或附点音符。
\end{solution}
